\section{Realización General del Proyecto}
    Se hace uso de la implementación en NumPy \parencite{numpy_library} de la 
    Transformada Rápida de Fourier (FFT) para generar el dominio de frecuencias 
    de cada una de las imágenes además de aplicarles el centrado del espectro. 
    Para determinar de mejor manera qué filtros de recorte serían lo más 
    convenientes en cada caso, se visualizaron las frecuencias para delimitar 
    las frecuencias dominantes o que introducían ruido periódico a las imágenes, 
    es decir, los puntos brillantes o destellos en la visualización.\\

    El filtro de Butterworth rechaza banda (BBRF) se implementó siguiendo tanto 
    las notas como el código proporcionado en clase \parencite{tinoco2026pdi}, el 
    cual se modificó ligeramente para evitar divisiones entre cero. Para crear 
    el filtro de recorte únicamente se compusieron los filtros de Butterworth con 
    diferente parámetros (centros, ancho de banda y orden), donde esta composición 
    es una multiplicación punto a punto de matrices que representan los filtros. \\

    Por último, para generar la imagen filtrada solamente se multiplica su dominio 
    de frecuencias por el filtro de recorte definido y se invierte la FFT usando 
    la propia implementación de NumPy. En la visualización se considera tanto el 
    espectro filtrado como la imagen resultante para compararla con la imagen original.